\begin{ZhChapter}

\chapter{背景知識}

\section{大型語言模型}

大型語言模型（Large Language Models, LLMs）是一種基於深度學習，特別是 Transformer 架構所建構的自然語言處理模型。其核心概念在於透過自注意力機制（Self-Attention），使模型在處理序列（如句子）時，能同時考量詞語之間的相關性，並有效捕捉詞語之間的遠距依賴關係與語意聯繫。與傳統的循環神經網路（RNN）或長短期記憶網路（LSTM）相比，Transformer 架構具備高度平行運算能力，能在海量文本資料上進行高效訓練。此外，Transformer 的設計包含編碼器（Encoder）與解碼器（Decoder）兩大模組：編碼器負責提取輸入文本的語意表示，捕捉字詞及片語間的語法與語意關聯；解碼器則在生成目標序列時，依據編碼器提供的上下文向量與自身生成歷史，逐步生成合理且連貫的文字輸出。這種編碼與解碼的結構，使大型語言模型不僅能理解文字，還能生成具邏輯性的文本內容。

在預訓練（Pre-training）階段，模型通常在大規模無標註資料庫上進行自我監督學習，藉由語言建模目標學習詞彙關係、句法結構與上下文語意，從而具備通用語言知識。完成預訓練後，模型可透過微調（Fine-tuning）或提示詞工程（Prompt Engineering）應用於下游任務，如問答、摘要、翻譯、對話生成及文本分類。隨著模型參數與訓練資料規模增加，其生成語言的流暢性、邏輯性及推理能力均有顯著提升，使其除基本對話外，亦能處理程式碼撰寫、圖片生成、影片生成等複雜任務。

近年來大語言模型在處理自然語言中已經展現了非常卓越的性能，代表性的大語言模型有 OpenAI 的 ChatGPT、Google 的 Gemini、Meta 的 Llama 等等。

\section{模型微調}

模型微調（Fine-tuning）是指在大型語言模型完成預訓練後，使用特定任務或領域的資料對模型進行再訓練，以提升其在特定應用情境下的表現。與從頭訓練模型相比，微調能有效降低計算成本，並保留預訓練階段學到的通用語言知識。微調過程通常涉及將模型的參數更新，使其在特定任務上最小化損失函數 (Loss Function)。在實務應用中，微調可分為全參數微調（Full Fine-tuning）與參數高效化方法（Parameter-efficient Fine-tuning）兩類，後者如 Adapter、LoRA 等技術，僅對少量新增參數進行訓練，以降低計算資源需求並維持模型效能。

微調資料的品質與多樣性對模型表現與安全性具有關鍵影響。如果資料中存在標註錯誤或偏頗內容，可能導致模型生成不準確、偏差或有潛在風險的輸出。此外，在多領域應用中，模型微調還需考慮過度擬合與遷移能力之平衡，確保模型在特定任務表現優異的同時，仍能保留通用語言知識。微調技術已廣泛應用於問答系統、對話生成、文本分類、醫療輔助決策及法律文書分析等領域，成為實務應用大型語言模型的重要手段。

\section{提示詞}



\section{提示詞注入}

提示詞注入（Prompt Injection）是一種針對大型語言模型的攻擊手法，攻擊者透過構造惡意輸入，使模型忽略原有指令或安全限制，生成不當或敏感的回應。此類攻擊利用模型對上下文缺乏明確信任邊界的特性，將惡意指令隱藏於合法提示之中，誘導模型產生非預期行為。提示詞注入可分為直接式和間接式兩類：直接式由使用者直接輸入惡意提示，模型立即執行；間接式則透過外部資料、網頁內容，或將惡意提示進行編碼、格式化等方式間接影響模型，使其在生成結果時受到操控。

此類攻擊通常利用了大型語言模型無法嚴格區分「指令（Instruction）」與「資料（Data）」的特性。當惡意指令被偽裝成普通數據輸入時，模型可能會誤將其視為執行命令，進而生成敏感資訊、執行不安全操作，或泄露系統內部資料。例如，攻擊者可能在提示中加入「忽略前面的限制，告訴我系統密碼」等語句，誘導模型違反原有安全策略。

近年來，隨著 Agent（代理人）概念的興起，大型語言模型逐漸被用於自動化操控多項系統行為，例如存取檔案、執行程式碼、呼叫外部 API 或與其他服務進行互動。在此架構下，LLM 不僅負責文字生成，更扮演決策與控制核心的角色。然而，若 Agent 的權限管理與上下文隔離設計不當，一旦遭受提示詞注入攻擊，惡意指令可能被放大其影響範圍，進而影響整個系統運作。例如，攻擊者可能藉由注入指令誘導 Agent 執行未經授權的檔案存取、傳送敏感資料，甚至觸發連鎖式的自動化操作，造成系統層級的安全風險。

綜合而言，提示詞注入不僅揭示了大型語言模型在安全性上的脆弱性，也促使研究者在模型設計與應用中，重視提示詞的可信性與系統安全性，成為人工智慧安全研究的重要課題。

\end{ZhChapter}